% Chapter 3: Comprehensive Validation Framework
% Double Machine Learning for Time Series - Volume 2

\chapter{Comprehensive Validation Framework}

% ============================================================================
% SECTION 1: INTRODUCTION TO VALIDATION STRATEGY
% ============================================================================

\section{Introduction to Validation Strategy}

Before applying Double Machine Learning (DML) to real-world empirical data---where the true treatment effect is unknown---we must establish confidence in our methodology through rigorous validation on synthetic data where the ground truth is known by construction. This chapter presents a comprehensive validation framework comparing seven causal inference methods, demonstrating why DML's cross-fitting and Neyman orthogonality provide superior performance under confounding.

\subsection{Why Validate Causal Inference Methods?}

Causal inference methods make strong assumptions: unconfoundedness, overlap, correct functional forms for nuisance parameters. Unlike supervised learning where prediction accuracy can be measured directly on held-out test sets, causal effects cannot be validated through simple train-test splits---the fundamental problem of causal inference is that we never observe both $Y_i(1)$ and $Y_i(0)$ for the same unit.

Validation on synthetic data addresses this challenge by constructing data generating processes (DGPs) where:
\begin{itemize}
    \item The true treatment effect $\tau_0$ is known by design
    \item Confounding strength is controllable
    \item Functional forms (linear vs nonlinear) can be varied
    \item Sample sizes and dimensions can be systematically tested
\end{itemize}

This allows us to evaluate estimator properties---bias, variance, coverage---that would be impossible to assess with real data alone.

\subsection{Synthetic vs Empirical Validation Approaches}

Our validation strategy follows a two-stage approach:

\textbf{Stage 1 (This Chapter): Synthetic Validation}

We generate controlled datasets with known treatment effects and systematically test estimators across scenarios:
\begin{itemize}
    \item \textbf{Confounding strength}: From none ($\gamma = 0$) to strong ($\gamma = 2$)
    \item \textbf{Sample sizes}: Small ($n=200$) to large ($n=5000$)
    \item \textbf{Functional forms}: Linear and nonlinear DGPs
    \item \textbf{Dimensionality}: Low ($p=5$) to moderate ($p=30$) confounders
\end{itemize}

\textbf{Stage 2 (Chapter 4): Empirical Replication}

After establishing DML's superior performance on synthetic data, we validate against published empirical benchmarks---specifically, Chernozhukov et al. (2018) 401(k) study---where expert consensus provides an approximate ``ground truth'' for comparison.

This synthetic-first approach builds reader confidence: if DML works correctly when we know $\tau_0$, we can trust its estimates when we don't.

\subsection{The Seven-Method Comparison Framework}

We compare DML against six baseline methods spanning parametric, semi-parametric, and nonparametric approaches. Table~\ref{tab:method_comparison} summarizes key properties.

\begin{table}[h]
\centering
\caption{Comparison of Seven Causal Inference Methods}
\label{tab:method_comparison}
\begin{tabular}{llcccc}
\hline
\textbf{Method} & \textbf{Type} & \textbf{Cross-Fit} & \textbf{Robust} & \textbf{Speed} & \textbf{Complexity} \\
\hline
NaiveOLS & Parametric & No & No & Fast & Low \\
OLSWithControls & Parametric & No & Partial & Fast & Low \\
IPW & Semi-parametric & No & Partial & Moderate & Moderate \\
AIPW & Semi-parametric & No & Yes* & Moderate & Moderate \\
RandomForest & Nonparametric & No & Partial & Slow & High \\
XGBoost & Nonparametric & No & Partial & Moderate & High \\
\textbf{DML} & \textbf{Semi-parametric} & \textbf{Yes} & \textbf{Yes} & \textbf{Moderate} & \textbf{Moderate} \\
\hline
\multicolumn{6}{l}{\small *Doubly robust without cross-fitting; DML adds sample-splitting for orthogonality}
\end{tabular}
\end{table}

\textbf{Parametric Methods}:
\begin{itemize}
    \item \textbf{NaiveOLS}: Simple regression $Y \sim T$ (ignores confounding)---serves as worst-case baseline
    \item \textbf{OLSWithControls}: Linear regression $Y \sim T + X$---standard econometric approach
\end{itemize}

\textbf{Semi-Parametric Methods}:
\begin{itemize}
    \item \textbf{IPW}: Inverse propensity weighting---reweights to balance treatment groups
    \item \textbf{AIPW}: Augmented IPW---doubly robust but without cross-fitting
\end{itemize}

\textbf{Nonparametric ML Baselines}:
\begin{itemize}
    \item \textbf{RandomForest}: Flexible tree ensemble without cross-fitting
    \item \textbf{XGBoost}: Gradient boosted trees without cross-fitting
\end{itemize}

\textbf{DML (Our Focus)}:
\begin{itemize}
    \item Semi-parametric framework with Neyman orthogonality
    \item Cross-fitting eliminates overfitting bias
    \item Flexible ML for nuisance parameters ($\mu_0, e_0$)
    \item Asymptotic normality even when $p \approx n^{1/4}$
\end{itemize}

The critical distinction: \textit{cross-fitting}. While AIPW achieves double robustness (correct inference if either outcome or propensity model is correct), DML's sample-splitting ensures nuisance parameter estimates are independent of the score used for causal inference---a property essential for handling high-dimensional ML models.

\subsection{Chapter Roadmap}

This chapter progresses from simple baselines to sophisticated DML variants:

\textbf{Section 2: Baseline Methods Comparison} (~15 pages)

Two comprehensive experiments demonstrating baseline method limitations:
\begin{itemize}
    \item \textbf{Experiment 3.1}: Confounding strength sensitivity (5 scenarios)
    \item \textbf{Experiment 3.2}: Sample size robustness ($n \in \{200, 500, 1000, 2000, 5000\}$)
\end{itemize}

Expected findings: NaiveOLS fails under any confounding; OLS requires correct linear specification; IPW/AIPW sensitive to propensity score estimation; ML baselines improve but lack theoretical guarantees.

\textbf{Section 3: DML Deep Dive} (~18 pages)

Three experiments exploring DML mechanics:
\begin{itemize}
    \item \textbf{Experiment 3.3}: Cross-fitting sensitivity (2-10 folds)
    \item \textbf{Experiment 3.4}: Nuisance model comparison (RF vs Lasso vs XGB)
    \item \textbf{Experiment 3.5}: Treatment heterogeneity (ATE vs CATE)
\end{itemize}

Demonstrates why cross-fitting matters, how to choose nuisance models, and when DML captures heterogeneous effects.

\textbf{Section 4: Statistical Testing Framework} (~12 pages)

Rigorous hypothesis testing methodology:
\begin{itemize}
    \item \textbf{Experiment 3.6}: Statistical power analysis (100-2000 simulations)
    \item Multiple testing correction (Bonferroni/Holm)
    \item PASS/WARNING/FAIL interpretation
\end{itemize}

Establishes statistical rigor beyond descriptive comparisons.

\textbf{Section 5: Computational Performance} (~10 pages)

Practical runtime analysis:
\begin{itemize}
    \item \textbf{Experiment 3.7}: Computational benchmarks (all 7 methods)
    \item Parallel execution on 64-core hardware
    \item Scalability ($n, p, $ simulations)
\end{itemize}

Shows DML's moderate computational cost is justified by theoretical guarantees.

\textbf{Section 6: Practical Recommendations} (~8 pages)

Decision framework for practitioners:
\begin{itemize}
    \item When to use which method?
    \item Validation checklist
    \item Common pitfalls and diagnostics
    \item Transition to empirical applications
\end{itemize}

By the end of this chapter, readers will understand not just \textit{that} DML works, but \textit{why} it outperforms alternatives and \textit{when} to apply it.

\subsection{Running Example: Synthetic Validation Data}

Throughout this chapter, we use the \texttt{DGPGenerator} class to create controlled synthetic datasets. Here's how to generate data with known treatment effect $\tau_0 = 2.0$ and moderate confounding:

\begin{minted}{python}
import numpy as np
from src.validation.dgp_generator import DGPGenerator

# Set seed for reproducibility
np.random.seed(42)

# Create data generating process
dgp = DGPGenerator(
    n=1000,                     # Sample size
    p=5,                        # Number of confounders
    true_effect=2.0,           # Known treatment effect
    confounding_strength=1.0,   # Moderate confounding
    treatment_model='linear',   # Linear propensity score
    outcome_model='linear',     # Linear outcome function
    noise_level=1.0,           # Standard error of noise
    random_state=42
)

# Generate single dataset
data = dgp.generate()

print(f"Outcome Y: {data.Y.shape}")           # (1000,)
print(f"Treatment T: {data.T.shape}")         # (1000,) binary
print(f"Confounders X: {data.X.shape}")       # (1000, 5)
print(f"True effect: {data.true_effect}")     # 2.0
\end{minted}

This DGP constructs data where:
\[
T_i \sim \text{Bernoulli}\left(\text{logit}^{-1}(\gamma X_i'\beta)\right), \quad
Y_i = \tau_0 T_i + X_i'\alpha + \epsilon_i
\]
where $\gamma$ controls confounding strength (here $\gamma=1.0$), and $\beta, \alpha$ are randomly generated coefficient vectors ensuring $X$ affects both treatment assignment and outcomes.

Key design feature: when $\gamma = 0$ (no confounding), even NaiveOLS recovers $\tau_0$ correctly; as $\gamma$ increases, confounding bias grows, and only methods properly adjusting for $X$ remain unbiased.

We can verify the data generating process:

\begin{minted}{python}
# Run 10,000 simulations to check average treatment effect
np.random.seed(42)
dgp_check = DGPGenerator(n=1000, p=5, true_effect=2.0,
                          confounding_strength=1.0, random_state=42)

effects = []
for _ in range(10000):
    data = dgp_check.generate()

    # Oracle estimator: knows true functional form
    from sklearn.linear_model import LinearRegression
    model = LinearRegression()
    model.fit(np.c_[data.T, data.X], data.Y)
    tau_hat = model.coef_[0]  # Treatment coefficient
    effects.append(tau_hat)

print(f"Oracle bias: {np.mean(effects) - 2.0:.4f}")
print(f"Oracle std: {np.std(effects):.4f}")
# Output: Oracle bias: -0.0003 (essentially unbiased)
#         Oracle std: 0.1012 (standard error)
\end{minted}

The oracle estimator (knowing true linear specification) achieves near-zero bias. Our validation experiments test whether DML matches oracle performance when functional forms are \textit{unknown} and potentially misspecified.

\textbf{Next Steps}: Section 2 begins systematic comparison of all seven methods, starting with confounding sensitivity analysis.

% Transition to Section 2
The stage is set. We have synthetic data with known $\tau_0$, a battery of seven estimation methods, and a rigorous statistical testing framework. Section 2 puts these methods to the test.

